\uuid{Je94}
\exo7id{469}
\titre{exo7 469}
\auteur{cousquer}
\organisation{exo7}
\datecreate{2003-10-01}
\isIndication{false}
\isCorrection{false}
\chapitre{Propriétés de R}
\sousChapitre{Maximum, minimum, borne supérieure}
\cpge{mpsi-mp2i.nombres-reels-suites.b, pcsi.nombres-reels-suites.a, ptsi.nombres-reels-suites.a, mpsi-mp2i.calcul-algebrique-trigonometrie.c, pcsi.calcul-algebrique-trigonometrie.c, ptsi.calcul-algebrique-trigonometrie.c, bcpst1.nombres-reels}
\module{Analyse}
\niveau{L1}
\difficulte{}

\contenu{
\question{
On considère l'ensemble des nombres de la forme $1+\frac{1}{n}$,
où $n$ décrit l'ensemble des entiers strictement positifs. Cet ensemble
est-il majoré ? Minoré ? A-t-il un plus petit élément ? 
Un plus grand élément ? Justifier vos réponses.
}
}
